Setting. By Fact \claimref{standard-model.imported-pair-systems-two-chirality-halves} the pair system is the Fock module $\Lambda^{\bullet}W$ of three modes, with $\Lambda^{k}W$ transforming under $U(3)$ as $1, 3, \bar{3}, 1$ for $k = 0, 1, 2, 3$. Theorem \claimref{standard-model.global-form-quotient-determined} gives the abelian factor of $S(U(2) \times U(3))$ as $(\lambda^{3}I_{2}, \lambda^{-2}I_{3})$ up to the $\mathbb{Z}_{6}$ quotient, with $Y = q/6$ for the weight $q$; since $\mu I_{2}$ acts on $F_{n}$ as $\mu^{n}$ and $\mu I_{3}$ on $\Lambda^{k}W$ as $\mu^{k}$, the weight on $F_{n} \otimes \Lambda^{k}W$ is $q = 3n - 2k$. A polynomial irreducible $F_{\lambda}$ is homogeneous of degree $n = \lambda_{1} + \lambda_{2}$, restricts to $SU(2)$ as spin $(\lambda_{1} - \lambda_{2})/2$, and every finite polynomial representation is a sum of them. The particle half keeps $(n, k)$ with $\pi n + k$ even. The four conditions are those of Fact \claimref{standard-model.imported-anomaly-coefficients}, read with left-handed Weyl fermions.

(1) Normalisation: the Fact sums $Y$ over weak doublets weighted by colour multiplicity and over colour triplets and antitriplets weighted by weak multiplicity; for a weak representation of spin $j$ the $[SU(2)]^{2}U(1)$ coefficient is proportional to its Dynkin index, and the weight $\tau_{j}$, twice the index, gives the doublet weight one, so the Fact's sum is the case where only doublets occur. Take $\pi = 1$. For $n$ even the surviving pair grades are $k = 0, 2$, with colour dimensions $1, 3$ and hypercharges $\tfrac{n}{2}$ and $\tfrac{3n - 4}{6}$, so the colour-weighted sum of $Y$ is $\tfrac{n}{2} + \tfrac{3n - 4}{2} = 2(n - 1)$; for $n$ odd they are $k = 1, 3$ with dimensions $3, 1$ and hypercharges $\tfrac{3n - 2}{6}$, $\tfrac{3n - 6}{6}$, and the sum is again $2(n - 1)$. Take $\pi = 0$: the surviving grades are $k = 0, 2$ for every $n$, and the same sum is $2(n - 1)$. Multiplying by $d$ gives the gravitational sum, and by $\tau$ the $[SU(2)]^{2}U(1)$ sum, at either parity. The $[SU(3)]^{2}U(1)$ sum takes the single coloured grade with its weak multiplicity $d$: for $\pi = 1$, $k = 2$ ($Y = \tfrac{3n-4}{6}$) for $n$ even and $k = 1$ ($Y = \tfrac{3n - 2}{6}$) for $n$ odd, giving $\tfrac{d}{6}(3n - 3 - (-1)^{n})$; for $\pi = 0$, $k = 2$ always, giving $\tfrac{d}{6}(3n - 4)$. For $\pi = 1$ and $\lambda = (a, a)$, $n = 2a$, $d = 1$, the cubic sum is $a^{3} + 3(\tfrac{6a - 4}{6})^{3}$. The script recomputes all of these by direct summation over the components, at both parities, for every $\lambda$ with $\lambda_{1} \le 8$ and compares.

(2) Write $m_{\lambda}$ for the multiplicity of $F_{\lambda}$ in $F$. In the $[SU(2)]^{2}U(1)$ sum $\sum_{\lambda} m_{\lambda}\,2\tau_{\lambda}(n_{\lambda} - 1)$, the same at both parities, every term is nonnegative: $\tau_{\lambda} = 0$ for weak singlets, $n - 1 = 0$ for $\lambda = (1, 0)$, and the only other $\lambda$ with $n \le 1$ is $(0, 0)$, a singlet. So the condition holds if and only if $m_{\lambda} = 0$ for every $\lambda$ with $\lambda_{1} > \lambda_{2}$ and $n \ge 2$. What remains is $x = m_{(0,0)}$, $y = m_{(1,0)}$ and $z_{a} = m_{(a,a)}$, $a \ge 1$. The gravitational condition reads $-2x + \sum_{a} 2(2a - 1)z_{a} = 0$, so $x = \sum_{a}(2a - 1)z_{a}$, at both parities.

Even slots, $\pi = 0$. The $[SU(3)]^{2}U(1)$ coefficients of $(0,0)$, $(1,0)$, $(a,a)$ are $-\tfrac23$, $-\tfrac13$, $\tfrac{6a - 4}{6}$, so the condition, times six, reads $-4x - 2y + \sum_{a}(6a - 4)z_{a} = 0$, that is $y = -\sum_{a} a\,z_{a}$. Multiplicities are nonnegative, so $y = 0$ and every $z_{a} = 0$, hence $x = 0$: no nonzero $F$ is anomaly-free at even slot parity.

Odd slots, $\pi = 1$. The coefficients are $-\tfrac23$, $\tfrac13$, $\tfrac{6a - 4}{6}$, and the condition, times six, reads $-4x + 2y + \sum_{a}(6a - 4)z_{a} = 0$, so $y = \sum_{a} a\,z_{a}$. The cubic coefficients of $(0,0)$ and $(1,0)$ are $-\tfrac89$ and $-\tfrac29$, so the cubic condition becomes $\sum_{a} z_{a}\,g(a) = 0$ with
\[
g(a) = a^{3} + 3\bigl(a - \tfrac23\bigr)^{3} - \tfrac89(2a - 1) - \tfrac29 a = 4a^{3} - 6a^{2} + 2a = 2a(2a - 1)(a - 1),
\]
which vanishes at $a = 1$ and is positive for $a \ge 2$. Hence $z_{a} = 0$ for $a \ge 2$, and $x = y = z_{1} =: m$: $F \cong m\,(F_{(0,0)} \oplus F_{(1,0)} \oplus F_{(1,1)}) = (\Lambda^{\bullet}\C^{2})^{\oplus m}$, graded. Conversely for this $F$ at $\pi = 1$ all four sums vanish, being $m$ times those of $\Lambda^{\bullet}\C^{2}$, which are $(-2 + 0 + 2)$, $0$, $(-\tfrac23 + \tfrac13 + \tfrac13)$ and $(-\tfrac89 - \tfrac29 + \tfrac{10}{9})$. The listed content is the particle half read off component by component: $(0,0)$, $(0,2)$, $(1,1)$, $(1,3)$, $(2,0)$, $(2,2)$. The $[SU(3)]^{3}$ coefficient, $-1 + 2 - 1$ per copy with $\bar{3}$ counted $-1$, vanishes too. The script checks the classification independently by exhausting all multiplicities $0 \le m_{\lambda} \le 2$ over the six $\lambda$ with $\lambda_{1} \le 2$ at both parities, finding exactly $m = 1, 2$ at $\pi = 1$ and nothing at $\pi = 0$.

(3) For $F = (\Lambda^{\bullet}\C^{2})^{\oplus m}$, $F_{0} \cong \C^{m}$ (trivial), $F_{2} \cong \Lambda^{2}\C^{2} \otimes \C^{m}$ and $F_{n} = 0$ for $n \ge 3$. The symmetric part $a^{\dagger}(u)a^{\dagger}(u') + a^{\dagger}(u')a^{\dagger}(u)$ is a $U(2)$-equivariant map $\mathrm{Sym}^{2}\C^{2} \otimes F_{n} \to F_{n+2}$. For $n \ge 1$ the target is zero. For $n = 0$ it is an equivariant map from $F_{(2,0)}\otimes\C^{m}$ to $F_{(1,1)}\otimes\C^{m}$, zero by Schur's lemma since $(2,0) \neq (1,1)$ (imported, cited from memory; the script confirms that the highest weight $(2,0)$ does not occur among the weights of $\Lambda^{2}\C^{2}$). So the creation maps anticommute, and $a^{\dagger}(u)^{2} = 0$.

(4) Green's ansatz writes an order-$p$ parafermionic creation operator as $f_{i}^{\dagger} = \sum_{\alpha = 1}^{p} c_{i}^{\alpha\dagger}$ with $c^{\alpha}$ the creation operators of the $\alpha$-th of $p$ commuting copies of the two-mode fermionic Fock space; the Fock space is spanned by the words in $f_{1}^{\dagger}, f_{2}^{\dagger}$ applied to the vacuum, and $U(2)$ acts on each copy alike. The script builds these operators exactly for $p = 1, \dots, 4$, computes the dimension of each weight space $(n_{1}, n_{2})$ of the span, decomposes the weight multiset into $U(2)$ irreducibles by removing highest weights, and finds the $F_{\lambda}$ with $\lambda_{1} \le p$ once each (dimensions $4, 10, 20, 35$). For $p \ge 2$ the term $F_{(2,0)}$ appears, a weak triplet with $n = 2$, and by (2) the $[SU(2)]^{2}U(1)$ condition fails at either parity; the script lists the four sums at $p = 2$, $\pi = 1$, as $20$, $12$, $\tfrac{14}{3}$, $\tfrac{254}{9}$. That $F_{(2,0)}$ occurs for every $p \ge 2$ follows because the word $f_{1}^{\dagger}f_{1}^{\dagger}$ applied to the vacuum is $2\sum_{\alpha<\beta}c_{1}^{\alpha\dagger}c_{1}^{\beta\dagger}$ applied to it, nonzero once $p \ge 2$, a state of weight $(2, 0)$, which no $F_{\lambda}$ with $\lambda_{1} \le 1$ carries.

Scope. The assumption that the composite space is $F \otimes \Lambda^{\bullet}W$, with one pair system, $U(2)$ acting polynomially on the slots, slot operators of one parity $\pi$ and chirality the tower grading $(-1)^{\pi n + k}$, is part of the statement, not derived; it is weaker than Hypothesis (S) (\claimref{standard-model.hyp-s-slot-statistics}) in leaving the slot statistics free, and (S) is the case $\pi = 1$, $F = \Lambda^{\bullet}\C^{2}$, $m = 1$. The multiplicity $m$ is not fixed by the anomaly conditions, and that the five modes form one Clifford system is not addressed.
